% GENERATED FILE. Edit canonical lessons, then run npm run generate:books.
% canonical-source-hash: fnv1a64:dd4a9273f2e3f554
% canonical-lessons: SA-C433-ayah, SA-C433-muladhanam, SA-C433-amsah, SA-C433-bhagah, SA-C433-pratibhuh, SA-R433-review-426-433

\chapter{Income and Shares}
\label{ch:sa-income-and-shares}
\hlchaptermodality{433}

\begin{chapteropening}
\textbf{By the end of this chapter:} \emph{I can talk about income, capital and fair shares.}

Complete the last lesson of chapter 433: I can talk about income, capital and fair shares.
\end{chapteropening}

\section[\texorpdfstring{\sk{आयः} (āyaḥ)}{आयः (āyaḥ)}]{\sk{आयः} (āyaḥ) — income}

\label{lesson:SA-C433-ayah}



\begin{quote}
Before the new one: say the Sanskrit for saves up, then the Sanskrit for savings.
\end{quote}

\subsection*{You'll want to know: \sk{आयः}}
\textbf{\sk{आयः}} — \emph{āyaḥ} — \textquotedblleft{}income\textquotedblright{}.

\begin{quote}
\sk{कृषकस्य} \sk{आयः} \sk{अल्पः} \sk{अस्ति।}
\end{quote}

\begin{quote}
\emph{kṛṣakasya āyaḥ alpaḥ asti.}
\end{quote}

\begin{quote}
The farmer's income is small.
\end{quote}

A masculine noun.

\begin{cousinweb}
\textbf{\sk{आयः}} (\emph{āyaḥ}) is built from \emph{ā}, \textquotedblleft{}towards\textquotedblright{}, and the root √i, \textquotedblleft{}to go\textquotedblright{}: what comes in; MW gives \emph{āya} as arrival and as income, revenue (MW).
\end{cousinweb}

\begin{grammarlens}[title={What you've built}]
One more word for getting around.
\end{grammarlens}

\subsection*{Your turn}
\begin{itemize}
\raggedright
  \item \emph{Say it:} \emph{āyaḥ}
  \item \emph{Say it:} \emph{āyaḥ}, once more
  \item \emph{Say it:} \emph{kṛṣakasya āyaḥ alpaḥ asti.} — The farmer's income is small.
  \item \emph{Recall:} say the Sanskrit for saves up, then the Sanskrit for savings, then say \emph{āyaḥ} again
\end{itemize}

\begin{tcolorbox}[breakable,colback=teal!4,colframe=teal!35!black,title={Before you move on}]
What does \sk{आयः} mean? (\textquotedblleft{}Income\textquotedblright{}.) Say it once more.
\end{tcolorbox}
\section[\texorpdfstring{\sk{मूलधनम्} (mūladhanam)}{मूलधनम् (mūladhanam)}]{\sk{मूलधनम्} (mūladhanam) — trading capital}

\label{lesson:SA-C433-muladhanam}



\begin{quote}
Before the new one: say the Sanskrit for savings, then the Sanskrit for income.
\end{quote}

\subsection*{You'll want to know: \sk{मूलधनम्}}
\textbf{\sk{मूलधनम्}} — \emph{mūladhanam} — \textquotedblleft{}trading capital\textquotedblright{}.

\begin{quote}
\sk{वणिक्} \sk{मूलधनेन} \sk{व्यापारं} \sk{करोति।}
\end{quote}

\begin{quote}
\emph{vaṇik mūladhanena vyāpāraṃ karoti.}
\end{quote}

\begin{quote}
The merchant trades with his capital.
\end{quote}

A neuter noun.

\begin{cousinweb}
\textbf{\sk{मूलधनम्}} (\emph{mūladhanam}) is built from \emph{mūla}, \textquotedblleft{}a root, the base\textquotedblright{}, and \emph{dhana}, \textquotedblleft{}wealth\textquotedblright{}; MW gives the compound as principal, capital (MW).
\end{cousinweb}

\begin{grammarlens}[title={What you've built}]
One more word for getting around.
\end{grammarlens}

\subsection*{Your turn}
\begin{itemize}
\raggedright
  \item \emph{Say it:} \emph{mūladhanam}
  \item \emph{Say it:} \emph{mūladhanam}, once more
  \item \emph{Say it:} \emph{vaṇik mūladhanena vyāpāraṃ karoti.} — The merchant trades with his capital.
  \item \emph{Recall:} say the Sanskrit for savings, then the Sanskrit for income, then say \emph{mūladhanam} again
\end{itemize}

\begin{tcolorbox}[breakable,colback=teal!4,colframe=teal!35!black,title={Before you move on}]
What does \sk{मूलधनम्} mean? (\textquotedblleft{}Trading capital\textquotedblright{}.) Say it once more.
\end{tcolorbox}
\section[\texorpdfstring{\sk{अंशः} (aṃśaḥ)}{अंशः (aṃśaḥ)}]{\sk{अंशः} (aṃśaḥ) — a portion; a part}

\label{lesson:SA-C433-amsah}



\begin{quote}
Before the new one: say the Sanskrit for income, then the Sanskrit for trading capital.
\end{quote}

\subsection*{You'll want to know: \sk{अंशः}}
\textbf{\sk{अंशः}} — \emph{aṃśaḥ} — \textquotedblleft{}a portion; a part\textquotedblright{}.

\begin{quote}
\sk{सः} \sk{धनस्य} \sk{अंशं} \sk{मित्राय} \sk{ददाति।}
\end{quote}

\begin{quote}
\emph{saḥ dhanasya aṃśaṃ mitrāya dadāti.}
\end{quote}

\begin{quote}
He gives a portion of the money to his friend.
\end{quote}

A masculine noun.

\begin{cousinweb}
\textbf{\sk{अंशः}} (\emph{aṃśaḥ}) probably comes from the root √aś, \textquotedblleft{}to obtain\textquotedblright{}, as MW gives it; MW gives \emph{aṃśa} as a share, a portion, a part (MW).
\end{cousinweb}

\begin{grammarlens}[title={What you've built}]
One more word for getting around.
\end{grammarlens}

\subsection*{Your turn}
\begin{itemize}
\raggedright
  \item \emph{Say it:} \emph{aṃśaḥ}
  \item \emph{Say it:} \emph{aṃśaḥ}, once more
  \item \emph{Say it:} \emph{saḥ dhanasya aṃśaṃ mitrāya dadāti.} — He gives a portion of the money to his friend.
  \item \emph{Recall:} say the Sanskrit for income, then the Sanskrit for trading capital, then say \emph{aṃśaḥ} again
\end{itemize}

\begin{tcolorbox}[breakable,colback=teal!4,colframe=teal!35!black,title={Before you move on}]
What does \sk{अंशः} mean? (\textquotedblleft{}A portion; a part\textquotedblright{}.) Say it once more.
\end{tcolorbox}
\section[\texorpdfstring{\sk{भागः} (bhāgaḥ)}{भागः (bhāgaḥ)}]{\sk{भागः} (bhāgaḥ) — a share}

\label{lesson:SA-C433-bhagah}



\begin{quote}
Before the new one: say the Sanskrit for trading capital, then the Sanskrit for a portion.
\end{quote}

\subsection*{You'll want to know: \sk{भागः}}
\textbf{\sk{भागः}} — \emph{bhāgaḥ} — \textquotedblleft{}a share\textquotedblright{}.

\begin{quote}
\sk{मम} \sk{भागः} \sk{अल्पः} \sk{अस्ति।}
\end{quote}

\begin{quote}
\emph{mama bhāgaḥ alpaḥ asti.}
\end{quote}

\begin{quote}
My share is small.
\end{quote}

A masculine noun.

\begin{cousinweb}
\textbf{\sk{भागः}} (\emph{bhāgaḥ}) comes from the root √bhaj, \textquotedblleft{}to share, to divide\textquotedblright{}, which also gives \emph{bhajati}, \textquotedblleft{}shares\textquotedblright{}; MW gives \emph{bhāga} as a part, a portion, a share (MW).
\end{cousinweb}

\begin{grammarlens}[title={What you've built}]
One more word for getting around.
\end{grammarlens}

\subsection*{Your turn}
\begin{itemize}
\raggedright
  \item \emph{Say it:} \emph{bhāgaḥ}
  \item \emph{Say it:} \emph{bhāgaḥ}, once more
  \item \emph{Say it:} \emph{mama bhāgaḥ alpaḥ asti.} — My share is small.
  \item \emph{Recall:} say the Sanskrit for trading capital, then the Sanskrit for a portion, then say \emph{bhāgaḥ} again
\end{itemize}

\begin{tcolorbox}[breakable,colback=teal!4,colframe=teal!35!black,title={Before you move on}]
What does \sk{भागः} mean? (\textquotedblleft{}A share\textquotedblright{}.) Say it once more.
\end{tcolorbox}
\section[\texorpdfstring{\sk{प्रतिभूः} (pratibhūḥ)}{प्रतिभूः (pratibhūḥ)}]{\sk{प्रतिभूः} (pratibhūḥ) — a guarantor}

\label{lesson:SA-C433-pratibhuh}



\begin{quote}
Before the new one: say the Sanskrit for a portion, then the Sanskrit for a share.
\end{quote}

\subsection*{You'll want to know: \sk{प्रतिभूः}}
\textbf{\sk{प्रतिभूः}} — \emph{pratibhūḥ} — \textquotedblleft{}a guarantor\textquotedblright{}.

\begin{quote}
\sk{मम} \sk{मित्रं} \sk{मम} \sk{प्रतिभूः} \sk{भवति।}
\end{quote}

\begin{quote}
\emph{mama mitraṃ mama pratibhūḥ bhavati.}
\end{quote}

\begin{quote}
My friend stands as my guarantor.
\end{quote}

A masculine noun.

\begin{cousinweb}
\textbf{\sk{प्रतिभूः}} (\emph{pratibhūḥ}) is built from \emph{prati}, \textquotedblleft{}in place of\textquotedblright{}, and the root √bhū, \textquotedblleft{}to be\textquotedblright{}, which also gives \emph{bhavati}, \textquotedblleft{}is\textquotedblright{}: one who stands in for another; MW gives \emph{pratibhū} as a surety, security, bail (MW).
\end{cousinweb}

\begin{grammarlens}[title={What you've built}]
One more word for getting around.
\end{grammarlens}

\subsection*{Your turn}
\begin{itemize}
\raggedright
  \item \emph{Say it:} \emph{pratibhūḥ}
  \item \emph{Say it:} \emph{pratibhūḥ}, once more
  \item \emph{Say it:} \emph{mama mitraṃ mama pratibhūḥ bhavati.} — My friend stands as my guarantor.
  \item \emph{Recall:} say the Sanskrit for a portion, then the Sanskrit for a share, then say \emph{pratibhūḥ} again
\end{itemize}

\begin{tcolorbox}[breakable,colback=teal!4,colframe=teal!35!black,title={Before you move on}]
What does \sk{प्रतिभूः} mean? (\textquotedblleft{}A guarantor\textquotedblright{}.) Say it once more.
\end{tcolorbox}
\section[(dialogue)]{Looking back: from Abroad to Income and Shares}

\label{lesson:SA-R433-review-426-433}



\begin{quote}
Say the Sanskrit for a foreign country and crowded.
\end{quote}

\subsection*{Your turn}
Say these aloud:

\begin{itemize}
\raggedright
  \item a foreign country, crowded, a customs house, stays on, a place of refuge
  \item a night's stay, a tent, an encampment, spreads out, wakes up
  \item the master of the house, a visitor, puts up for the night, goes out to meet, sees a guest off
  \item unlocks, a clock, a travellers' lodge, a door bolt, rent
  \item a bazaar, a seller, a buyer, a customer, haggles
  \item worth, costly, low-priced, mixes, a touchstone
  \item an expense, wasteful spending, spends, saves up, savings
  \item income, trading capital, a portion, a share, a guarantor
\end{itemize}

\begin{tcolorbox}[breakable,colback=teal!4,colframe=teal!35!black,title={Before you move on}]
A foreign country? (\textbf{\sk{विदेशः}}, \emph{videśaḥ}.) Crowded? (\textbf{\sk{जनाकीर्णः}}, \emph{janākīrṇaḥ}.) A customs house? (\textbf{\sk{शुल्कस्थानम्}}, \emph{śulkasthānam}.) Stays on? (\textbf{\sk{अवतिष्ठते}}, \emph{avatiṣṭhate}.) A place of refuge? (\textbf{\sk{प्रतिश्रयः}}, \emph{pratiśrayaḥ}.) A night's stay? (\textbf{\sk{वसतिः}}, \emph{vasatiḥ}.) A tent? (\textbf{\sk{पटमण्डपः}}, \emph{paṭamaṇḍapaḥ}.) An encampment? (\textbf{\sk{निवेशः}}, \emph{niveśaḥ}.) Spreads out? (\textbf{\sk{आस्तृणाति}}, \emph{āstṛṇāti}.) Wakes up? (\textbf{\sk{प्रबुध्यते}}, \emph{prabudhyate}.) The master of the house? (\textbf{\sk{गृहपतिः}}, \emph{gṛhapatiḥ}.) A visitor? (\textbf{\sk{अभ्यागतः}}, \emph{abhyāgataḥ}.) Puts up for the night? (\textbf{\sk{निवासयति}}, \emph{nivāsayati}.) Goes out to meet? (\textbf{\sk{प्रत्युद्गच्छति}}, \emph{pratyudgacchati}.) Sees a guest off? (\textbf{\sk{अनुव्रजति}}, \emph{anuvrajati}.) Unlocks? (\textbf{\sk{उद्घाटयति}}, \emph{udghāṭayati}.) A clock? (\textbf{\sk{घटिका}}, \emph{ghaṭikā}.) A travellers' lodge? (\textbf{\sk{आवसथः}}, \emph{āvasathaḥ}.) A door bolt? (\textbf{\sk{अर्गलः}}, \emph{argalaḥ}.) Rent? (\textbf{\sk{अवक्रयः}}, \emph{avakrayaḥ}.) A bazaar? (\textbf{\sk{हट्टः}}, \emph{haṭṭaḥ}.) A seller? (\textbf{\sk{विक्रेता}}, \emph{vikretā}.) A buyer? (\textbf{\sk{क्रेता}}, \emph{kretā}.) A customer? (\textbf{\sk{ग्राहकः}}, \emph{grāhakaḥ}.) Haggles? (\textbf{\sk{पणते}}, \emph{paṇate}.) Worth? (\textbf{\sk{अर्घः}}, \emph{arghaḥ}.) Costly? (\textbf{\sk{महार्घः}}, \emph{mahārghaḥ}.) Low-priced? (\textbf{\sk{समर्घः}}, \emph{samarghaḥ}.) Mixes? (\textbf{\sk{मिश्रयति}}, \emph{miśrayati}.) A touchstone? (\textbf{\sk{निकषः}}, \emph{nikaṣaḥ}.) An expense? (\textbf{\sk{व्ययः}}, \emph{vyayaḥ}.) Wasteful spending? (\textbf{\sk{अपव्ययः}}, \emph{apavyayaḥ}.) Spends? (\textbf{\sk{व्ययति}}, \emph{vyayati}.) Saves up? (\textbf{\sk{सञ्चिनोति}}, \emph{sañcinoti}.) Savings? (\textbf{\sk{सञ्चयः}}, \emph{sañcayaḥ}.) Income? (\textbf{\sk{आयः}}, \emph{āyaḥ}.) Trading capital? (\textbf{\sk{मूलधनम्}}, \emph{mūladhanam}.) A portion? (\textbf{\sk{अंशः}}, \emph{aṃśaḥ}.) A share? (\textbf{\sk{भागः}}, \emph{bhāgaḥ}.) A guarantor? (\textbf{\sk{प्रतिभूः}}, \emph{pratibhūḥ}.)
\end{tcolorbox}
